\documentclass[11pt]{article}
\usepackage[margin=1in]{geometry}
\usepackage[T1]{fontenc}
\usepackage[utf8]{inputenc}
\usepackage{amsmath}
\usepackage{amssymb}
\usepackage{array}
\usepackage{booktabs}
\usepackage{hyperref}
\usepackage{longtable}
\usepackage{enumitem}

\title{Canonical Physics and Conventions Reference for \texttt{cqed\_sim}}
\author{Post-cleanup library reference}
\date{March 10, 2026}

\begin{document}
\maketitle

\noindent This report is the source of truth for the cleaned \texttt{cqed\_sim} library. It describes the conventions implemented in \texttt{cqed\_sim/}, distinguishes them from example-side code now under \texttt{examples/}, and records which conventions are directly backed by tests.

\section{Overview of cqed\_sim}
\texttt{cqed\_sim/} now contains only reusable simulator code:
\begin{itemize}[leftmargin=1.5em]
  \item \texttt{core}: Hilbert-space conventions, frames, Hamiltonian model, ideal gates, manifold-frequency helpers.
  \item \texttt{pulses}, \texttt{sequence}, \texttt{sim}: pulse objects, envelopes, builders, compilation, hardware distortion, solver path, extractors, noise.
  \item \texttt{calibration}, \texttt{observables}, \texttt{operators}, \texttt{tomo}, \texttt{io}, \texttt{plotting}, \texttt{unitary\_synthesis}: reusable helpers that still belong to the library surface.
\end{itemize}
Workflow, audit, study, paper-reproduction, and smoke-test code was moved to \texttt{examples/}, notably:
\texttt{examples/workflows/}, \texttt{examples/audits/}, \texttt{examples/studies/}, \texttt{examples/paper\_reproductions/}, and \texttt{examples/smoke\_tests/}.

\section{Hilbert Space Structure}
\begin{itemize}[leftmargin=1.5em]
  \item Runtime ordering is qubit first, cavity second:
  \[
  |q,n\rangle = |q\rangle_{\mathrm{qubit}} \otimes |n\rangle_{\mathrm{cavity}}.
  \]
  Implemented in \texttt{cqed\_sim/core/conventions.py} and \texttt{cqed\_sim/core/model.py}.
  \item Flattened basis index is \(\texttt{index}(q,n)=q\,n_{\mathrm{cav}}+n\). The pair \(\{|g,n\rangle,|e,n\rangle\}\) occupies \((n, n_{\mathrm{cav}}+n)\).
  \item Computational basis is \(|g\rangle=|0\rangle\), \(|e\rangle=|1\rangle\).
  \item The synthesis layer keeps the same full-space ordering but uses a restricted subspace basis ordered as \(|g,0\rangle,|e,0\rangle,|g,1\rangle,|e,1\rangle,\ldots\) in \texttt{cqed\_sim/unitary\_synthesis/subspace.py}.
\end{itemize}
Direct tests: \texttt{tests/test\_25\_tensor\_product\_convention.py}, \texttt{tests/unitary\_synthesis/test\_subspace.py}, \texttt{examples/audits/tests/test\_experiment\_convention\_audit.py}.

\section{Operator Conventions}
\begin{itemize}[leftmargin=1.5em]
  \item The simulator uses QuTiP Pauli matrices without sign changes. Therefore
  \[
  \sigma_z|g\rangle=+|g\rangle,\qquad \sigma_z|e\rangle=-|e\rangle.
  \]
  \item The core Hamiltonian uses the excitation projector \(n_q=b^\dagger b\), not \(\sigma_z/2\). For two levels,
  \[
  n_q = |e\rangle\langle e| = \frac{I-\sigma_z}{2}.
  \]
  \item Embedded drive operators are \(a\) for the cavity and \(b\) for the qubit, with runtime couplings \(a^\dagger \epsilon + a\epsilon^*\) and \(b^\dagger \epsilon + b\epsilon^*\).
  \item Ideal \(R_{xy}\), SQR, SNAP, and displacement gates are defined in \texttt{cqed\_sim/core/ideal\_gates.py}.
\end{itemize}
Direct tests: \texttt{tests/test\_16\_ideal\_primitives\_and\_extractors.py}, \texttt{tests/test\_20\_gaussian\_iq\_convention.py}, \texttt{tests/test\_25\_tensor\_product\_convention.py}.

\section{Hamiltonian Definitions}
The reusable runtime Hamiltonian in \texttt{cqed\_sim/core/model.py} is
\[
H_0 = \delta_c n_c + \delta_q n_q + \frac{\alpha}{2} b^{\dagger 2}b^2 + K(n_c) - \Xi(n_c)n_q,
\]
with
\[
K(n_c)= \frac{\texttt{kerr}}{2}n_c(n_c-1)+\cdots,\qquad
\Xi(n_c)= \texttt{chi}\,n_c + \texttt{chi2}\,n_c(n_c-1)+\texttt{chi3}\,n_c(n_c-1)(n_c-2)+\cdots .
\]
\texttt{cqed\_sim/core/frequencies.py} makes the intended transition law explicit:
\[
\omega_{ge}(n)=\omega_{ge}(0)-\Xi(n).
\]

\subsection*{Textbook-to-code translation}
With the active runtime sign convention,
\[
n_q=\frac{I-\sigma_z}{2}.
\]
Therefore
\[
\delta_q n_q = \frac{\delta_q}{2}I - \frac{\delta_q}{2}\sigma_z,\qquad
-\Xi(n_c)n_q = -\frac{\Xi(n_c)}{2}I + \frac{\Xi(n_c)}{2}\sigma_z.
\]
So, up to identity and cavity-only shifts,
\[
H_{q,\mathrm{runtime}} \doteq -\frac{1}{2}\bigl(\delta_q-\Xi(n_c)\bigr)\sigma_z.
\]
If one rewrites the same model with the textbook sign
\[
\sigma_z^{\mathrm{text}}=|e\rangle\langle e|-|g\rangle\langle g|=-\sigma_z,
\]
then
\[
H_{q,\mathrm{runtime}} \doteq \frac{1}{2}\bigl(\delta_q-\Xi(n_c)\bigr)\sigma_z^{\mathrm{text}}.
\]
Hence runtime \texttt{chi} is the per-photon downward pull of the \(|g,n\rangle \leftrightarrow |e,n\rangle\) transition.

\subsection*{Comparison with the synthesis drift convention}
\texttt{cqed\_sim/unitary\_synthesis/sequence.py} defines
\[
E_{g,n}= \delta_c n-\frac{\delta_q}{2}-\Xi_{\mathrm{synth}}(n)+K(n),\qquad
E_{e,n}= \delta_c n+\frac{\delta_q}{2}+\Xi_{\mathrm{synth}}(n)+K(n),
\]
and its docstring uses the opposite qubit sign convention:
\[
\sigma_z^{\mathrm{synth}}|e\rangle=+|e\rangle,\qquad \sigma_z^{\mathrm{synth}}|g\rangle=-|g\rangle.
\]
Its transition shift is therefore
\[
E_{e,n}-E_{g,n}=\delta_q+2\Xi_{\mathrm{synth}}(n).
\]
To match the runtime shift \(\delta_q-\Xi_{\mathrm{runtime}}(n)\), the algebraic mapping is
\[
2\Xi_{\mathrm{synth}}(n)=-\Xi_{\mathrm{runtime}}(n),
\]
so, componentwise,
\[
\chi_{\mathrm{synth}}=-\chi_{\mathrm{runtime}}/2,\qquad
\chi_{2,\mathrm{synth}}=-\chi_{2,\mathrm{runtime}}/2,\qquad
\chi_{3,\mathrm{synth}}=-\chi_{3,\mathrm{runtime}}/2.
\]
This mapping is derived from code; there is no dedicated runtime-to-synthesis translation helper yet.

\subsection*{Canonical parameter-definition table}
\small
\begin{longtable}{p{0.12\linewidth} p{0.24\linewidth} p{0.28\linewidth} p{0.24\linewidth}}
\toprule
Name & Code-level meaning & Algebra used & Caveat \\
\midrule
\endfirsthead
\toprule
Name & Code-level meaning & Algebra used & Caveat \\
\midrule
\endhead
\texttt{chi} & first-order runtime dispersive pull & \(-\texttt{chi}\,n_c n_q\) and \(\omega_{ge}(n)=\omega_{ge}(0)-\chi n\) & projector form, not stored as \(\sigma_z a^\dagger a/2\) \\
\texttt{chi2} & second-order pull, exposed through \texttt{chi\_higher[0]} & \(-\chi_2 n_c(n_c-1)n_q\) & falling-factorial coefficient \\
\texttt{chi3} & third-order pull, exposed through \texttt{chi\_higher[1]} & \(-\chi_3 n_c(n_c-1)(n_c-2)n_q\) & falling-factorial coefficient \\
\texttt{kerr} & first cavity Kerr coefficient & \(\texttt{kerr}\,n_c(n_c-1)/2\) & no implicit sign flip \\
higher Kerr & tuple \texttt{kerr\_higher} & \(K_m n_c^{\underline{m}}/m!\) & falling-factorial hierarchy \\
\texttt{alpha} & transmon Duffing anharmonicity & \(\texttt{alpha}\,b^{\dagger 2}b^2/2\) & inactive for \texttt{n\_tr=2} \\
\(\sigma_z\) & active runtime Pauli-\(Z\) & \(|g\rangle\langle g|-|e\rangle\langle e|=I-2n_q\) & opposite sign from synthesis docstring \\
\(n_q\) & qubit excitation projector & \(b^\dagger b\), and for two levels \(n_q=|e\rangle\langle e|\) & repo-specific versus \(\sigma_z/2\) form \\
\(\delta_q\) & qubit residual detuning & \(\omega_q-\omega_q^{\mathrm{frame}}\) & only defined after choosing a frame \\
\(\delta_c\) & cavity residual detuning & \(\omega_c-\omega_c^{\mathrm{frame}}\) & same frame caveat \\
\bottomrule
\end{longtable}
\normalsize

\subsection*{Standard-compatible vs repo-specific}
\begin{itemize}[leftmargin=1.5em]
  \item Standard-compatible: tensor ordering, QuTiP Pauli matrices, rad/s and seconds internally, Lindblad form, Duffing/Kerr terms, and \texttt{chi} as a transition pull.
  \item Repo-specific: projector-form \(n_q\) Hamiltonian, falling-factorial storage for higher-order pulls, and the explicit complex-envelope sign \(e^{+i(\omega t+\phi)}\).
  \item Needs caution: the synthesis drift model does not share the runtime \(\sigma_z\) sign or \(\chi\) normalization.
\end{itemize}

\section{Frequency and Time Units}
\begin{itemize}[leftmargin=1.5em]
  \item Internal Hamiltonian coefficients and frame frequencies are rad/s.
  \item Pulse durations, timeline steps, and solver \texttt{max\_step} are seconds.
  \item Example/workflow config dictionaries often use \texttt{\_hz} and \texttt{\_ns} names, then convert to core units before constructing the core model.
\end{itemize}
Direct tests: \texttt{tests/test\_05\_detuning\_and\_frames.py}, \texttt{tests/test\_10\_chi\_convention.py}, \texttt{tests/test\_13\_hardware\_extended.py}.

\section{Rotating Frame Conventions}
\begin{itemize}[leftmargin=1.5em]
  \item Frames are represented by \texttt{FrameSpec(omega\_c\_frame, omega\_q\_frame)} in \texttt{cqed\_sim/core/frame.py}.
  \item The core Hamiltonian uses \(\delta_c=\omega_c-\omega_c^{\mathrm{frame}}\) and \(\delta_q=\omega_q-\omega_q^{\mathrm{frame}}\).
  \item \texttt{manifold\_transition\_frequency(...)} returns the qubit transition in the chosen frame.
  \item \texttt{build\_sqr\_multitone\_pulse(...)} defaults to the rotating frame \((\omega_c,\omega_q)\) when \texttt{use\_rotating\_frame=True}.
\end{itemize}
Direct tests: \texttt{tests/test\_05\_detuning\_and\_frames.py}, \texttt{tests/test\_21\_qubox\_convention\_reconciliation.py}.

\section{Drive Hamiltonians}
\begin{itemize}[leftmargin=1.5em]
  \item \texttt{cqed\_sim/pulses/pulse.py} samples analytic or discrete envelopes as
  \[
  \epsilon(t)=\texttt{amp}\,\texttt{envelope}(t)\,e^{i(\texttt{carrier}t+\texttt{phase})}.
  \]
  \item \texttt{cqed\_sim/sim/runner.py} inserts \(\epsilon(t)\) and \(\epsilon^*(t)\) into creation and annihilation terms.
  \item Gaussian qubit rotations use \(\theta=2\int \Omega(t)\,dt\).
  \item Square cavity displacements use the analytic mapping \(\alpha=-i\int \epsilon(t)\,dt\).
  \item SQR multitone pulses use
  \[
  \epsilon_{\mathrm{SQR}}(t)=g(t)\sum_n A_n e^{i\phi_n}e^{i\omega_n t},
  \qquad
  A_n=\theta_n/(2T)+\lambda_0 d\lambda_n.
  \]
\end{itemize}
Direct tests: \texttt{tests/test\_20\_gaussian\_iq\_convention.py}, \texttt{tests/test\_23\_sqr\_additive\_amplitude\_correction.py}, \texttt{examples/audits/tests/test\_experiment\_convention\_audit.py}.

\section{Dissipation and Lindblad Terms}
\begin{itemize}[leftmargin=1.5em]
  \item \texttt{NoiseSpec} in \texttt{cqed\_sim/sim/noise.py} uses \(T_1\), \(T_\phi\), \(\kappa\), and \(n_{\mathrm{th}}\).
  \item The implemented convention is \(\gamma_1=1/T_1\) and \(\gamma_\phi=1/(2T_\phi)\).
  \item Two-level dephasing uses \(\sqrt{\gamma_\phi}\,\sigma_z\) with \(\sigma_z=I-2n_q\); multi-level dephasing uses \(\sqrt{\gamma_\phi}\,n_q\).
  \item \texttt{simulate\_sequence(...)} uses \texttt{mesolve} when collapse operators are present or when the initial state is already a density matrix; otherwise it uses \texttt{sesolve}.
\end{itemize}
Direct tests: \texttt{tests/test\_14\_dissipation.py}, \texttt{tests/test\_18\_noise\_and\_sideband.py}.

\section{State Representation}
\begin{itemize}[leftmargin=1.5em]
  \item Runtime inputs may be kets or density matrices.
  \item The result wrapper stores the final state, optional history, expectation traces, and the raw QuTiP solver result.
  \item Extractors and observables assume subsystem \(0\) is the qubit and subsystem \(1\) is the cavity, matching the tensor convention.
\end{itemize}
Direct tests: \texttt{tests/test\_14\_dissipation.py}, \texttt{tests/test\_16\_ideal\_primitives\_and\_extractors.py}, \texttt{tests/test\_19\_fock\_tomo.py}, \texttt{examples/smoke\_tests/tests/test\_sanity.py}.

\section{Hilbert Space Truncation}
\begin{itemize}[leftmargin=1.5em]
  \item \texttt{n\_cav} truncates the cavity Fock space and \texttt{n\_tr} truncates the transmon space.
  \item SQR \(\theta\) and \(\phi\) arrays are padded to the cavity dimension in \texttt{cqed\_sim/pulses/calibration.py}.
  \item The synthesis layer distinguishes full Hilbert dimension from restricted subspace dimension in \texttt{cqed\_sim/unitary\_synthesis/subspace.py}.
\end{itemize}
Direct tests: \texttt{tests/test\_11\_model\_invariants.py}, \texttt{tests/test\_24\_displacement\_selective\_spectroscopy.py}, \texttt{tests/unitary\_synthesis/test\_subspace.py}.

\section{Test Suite Audit}
\small
\begin{longtable}{p{0.31\linewidth} p{0.54\linewidth} p{0.10\linewidth}}
\toprule
File & Targets, type, and main assertion & Status \\
\midrule
\endfirsthead
\toprule
File & Targets, type, and main assertion & Status \\
\midrule
\endhead
\texttt{tests/test\_01\_sanity\_and\_free.py} & \texttt{core/model.py}, \texttt{sim/runner.py}; Hermiticity, free evolution, Kerr analytic phase. & Passed \\
\texttt{tests/test\_02\_cavity\_drive\_and\_kerr.py} & cavity drive and Kerr numerical sanity. & Passed \\
\texttt{tests/test\_03\_dispersive\_and\_ramsey.py} & dispersive phase scaling and Ramsey pull. & Passed \\
\texttt{tests/test\_04\_xy\_phase\_and\_overlap.py} & XY phase and overlap cancellation conventions. & Passed \\
\texttt{tests/test\_05\_detuning\_and\_frames.py} & detuning sign and rotating-frame behavior. & Passed \\
\texttt{tests/test\_06\_leakage\_drag\_and\_higher\_order.py} & multilevel leakage, DRAG, higher-order-term smoke test. & Passed \\
\texttt{tests/test\_07\_convergence\_regression.py} & timestep convergence plus golden regression artifact. & Passed \\
\texttt{tests/test\_08\_timeline\_and\_hardware.py} & overlap handling and hardware distortion basics. & Passed \\
\texttt{tests/test\_09\_runtime\_policy.py} & runtime-budget fast-path regression. & Passed \\
\texttt{tests/test\_10\_chi\_convention.py} & direct runtime \(\chi\) sign and meaning. & Passed \\
\texttt{tests/test\_11\_model\_invariants.py} & model dimensions and invariants. & Passed \\
\texttt{tests/test\_12\_pulse\_semantics.py} & pulse area/shape semantics. & Passed \\
\texttt{tests/test\_13\_hardware\_extended.py} & ZOH, low-pass, sidebands, quantization, crosstalk. & Passed \\
\texttt{tests/test\_14\_dissipation.py} & \(T_1\), \(T_\phi\), cavity damping. & Passed \\
\texttt{tests/test\_15\_numerics\_performance.py} & max-step, compile cache, scaling smoke. & Passed \\
\texttt{tests/test\_16\_ideal\_primitives\_and\_extractors.py} & ideal displacement/SNAP/SQR plus extractors. & Passed \\
\texttt{tests/test\_17\_gate\_library\_integration.py} & gate loading and library integration. & Passed \\
\texttt{tests/test\_18\_noise\_and\_sideband.py} & sideband exchange and noise composition. & Passed \\
\texttt{tests/test\_19\_fock\_tomo.py} & tomography and leakage-calibration helpers. & Passed \\
\texttt{tests/test\_20\_gaussian\_iq\_convention.py} & Gaussian I/Q axis, phase, detuning sign. & Passed \\
\texttt{tests/test\_21\_qubox\_convention\_reconciliation.py} & runtime-vs-documented external convention reconciliation. & Passed \\
\texttt{tests/test\_22\_support\_aware\_objective.py} & support-aware objective wiring. & Passed \\
\texttt{tests/test\_23\_sqr\_additive\_amplitude\_correction.py} & additive \(d\lambda\) amplitude convention. & Passed \\
\texttt{tests/test\_24\_displacement\_selective\_spectroscopy.py} & selective-spectroscopy frequencies. & Passed \\
\texttt{tests/test\_25\_tensor\_product\_convention.py} & tensor ordering, block indices, \(\sigma_z\) sign. & Passed \\
\texttt{tests/unitary\_synthesis/test\_subspace.py} & restricted-basis ordering and round-trip embedding. & Passed \\
\texttt{tests/unitary\_synthesis/test\_metrics.py} & fidelity gauge choices and leakage metrics. & Passed \\
\texttt{tests/unitary\_synthesis/test\_primitives\_and\_backends.py} & synthesis primitives and ideal/pulse backends. & Passed \\
\texttt{tests/unitary\_synthesis/test\_optimization.py} & duration mapping and time optimization. & Passed \\
\texttt{tests/unitary\_synthesis/test\_synthesis\_and\_targets.py} & easy targets, leakage penalty, MPS targets. & Passed \\
\texttt{tests/unitary\_synthesis/test\_free\_evolve\_condphase.py} & free-evolution conditional phase formulas. & Passed \\
\texttt{tests/unitary\_synthesis/test\_hardware\_grid\_parallel.py} & time grid, slew, tone spacing, parallel multistart. & Passed \\
\texttt{tests/unitary\_synthesis/test\_progress\_reporting.py} & history aggregation and serialization. & Passed \\
\texttt{tests/unitary\_synthesis/test\_reporting\_and\_determinism.py} & report schema and fixed-seed determinism. & Passed \\
\texttt{tests/unitary\_synthesis/test\_time\_policy\_and\_condphase.py} & time-policy modes and drift-phase decomposition. & Passed \\
\texttt{examples/audits/tests/test\_experiment\_convention\_audit.py} & moved audit checks for SU(2), tensor order, waveform signs, relative phases. & Passed \\
\texttt{examples/workflows/tests/test\_sqr\_transfer\_artifact.py} & moved workflow artifact serialization and replay. & Passed \\
\texttt{examples/studies/tests/test\_snap\_opt\_converges.py} & moved study regression for convergence from vanilla initialization. & Passed \\
\texttt{examples/studies/tests/test\_snap\_opt\_achieves\_zero\_coherent\_error\_above\_limit.py} & moved study regression for threshold behavior above the empirical limit. & Passed \\
\texttt{examples/studies/tests/test\_snap\_opt\_generalization\_sweeps.py} & moved study regressions for below-limit failure, runtime budget, and parameter sweeps. & Passed \\
\texttt{examples/studies/tests/test\_snap\_opt\_preserves\_low\_frequency\_content.py} & moved study convention test that no extra tones are introduced. & Passed \\
\texttt{examples/paper\_reproductions/tests/test\_convergence\_and\_runtime.py} & moved PRL-133 convergence trend and runtime-budget regression. & Passed \\
\texttt{examples/paper\_reproductions/tests/test\_no\_extra\_frequencies.py} & moved PRL-133 tone-support convention test. & Passed \\
\texttt{examples/paper\_reproductions/tests/test\_optimization\_limit.py} & moved PRL-133 optimization-limit regression. & Passed \\
\texttt{examples/paper\_reproductions/tests/test\_prl133\_components\_consistency.py} & moved PRL-133 component-consistency regression. & Passed \\
\texttt{examples/paper\_reproductions/tests/test\_prl133\_metric\_bounds.py} & moved PRL-133 metric boundedness and fail-fast validation. & Passed \\
\texttt{examples/paper\_reproductions/tests/test\_prl133\_metric\_relation.py} & moved PRL-133 metric-relation regression. & Passed \\
\texttt{examples/paper\_reproductions/tests/test\_prl133\_unitarity\_trace.py} & moved PRL-133 closed-system unitarity and trace-preservation check. & Passed \\
\texttt{examples/paper\_reproductions/tests/test\_reproduce\_baseline.py} & moved PRL-133 end-to-end reproduction smoke test. & Passed \\
\texttt{examples/smoke\_tests/tests/test\_sanity.py} & moved sequential-workflow sanity checks. & Passed \\
\texttt{examples/smoke\_tests/tests/test\_sqr\_calibration.py} & moved SQR calibration smoke and guarded-benchmark checks. & Passed \\
\texttt{examples/smoke\_tests/tests/test\_sqr\_speedlimit\_analysis.py} & moved speed-limit serial/parallel equivalence check. & Passed \\
\bottomrule
\end{longtable}
\normalsize

\section{Test Execution Results}
Commands used:
\begin{verbatim}
E:\Programs\python.exe -m pytest tests -q
E:\Programs\python.exe -m pytest examples/audits/tests examples/workflows/tests \
  examples/studies/tests examples/paper_reproductions/tests \
  examples/smoke_tests/tests -q
\end{verbatim}
Observed results:
\begin{itemize}[leftmargin=1.5em]
  \item \texttt{tests/}: 157 collected, 156 passed, 1 skipped, 0 failed.
  \item \texttt{examples/*/tests}: 49 collected, 49 passed, 0 failed.
  \item The single skip is the optional external MPS-reference check in \texttt{tests/unitary\_synthesis/test\_synthesis\_and\_targets.py}.
  \item Warning-only notes: two \texttt{numpy.trapz} deprecation warnings in \texttt{tests/test\_12\_pulse\_semantics.py} and \texttt{tests/test\_20\_gaussian\_iq\_convention.py}.
\end{itemize}
Interpreter note: \verb|E:\Program Files\Python311\python.exe| from \texttt{AGENTS.md} is not present in this workspace; tests were run with \verb|E:\Programs\python.exe| (Python 3.11.8).

\section{Verified Behaviors Backed by Tests}
\begin{itemize}[leftmargin=1.5em]
  \item Tensor ordering, basis labels, and synthesis restricted-basis ordering are directly tested.
  \item The active Pauli/Bloch sign convention is directly tested.
  \item Runtime \(\chi\) meaning as a downward transition pull is directly tested.
  \item The complex-envelope sign \(e^{+i(\omega t+\phi)}\) and Gaussian I/Q conventions are directly tested.
  \item The additive SQR amplitude correction rule is directly tested.
  \item Rotating-frame detuning behavior is directly tested.
  \item Solver selection and Lindblad behavior are directly tested.
  \item The moved example workflows still execute correctly in their new locations.
\end{itemize}

\section{Gaps Between Code Conventions and Test Coverage}
\begin{itemize}[leftmargin=1.5em]
  \item No dedicated test yet enforces the algebraic mapping between runtime \(\chi\) and synthesis-drift \(\chi\).
  \item Higher-order \(\chi_2\), \(\chi_3\), and higher-order Kerr coefficients have coverage, but not with isolated analytic sign/factor tests.
  \item Multi-level pure-dephasing behavior for \(\texttt{n\_tr}>2\) is only indirectly covered.
  \item The projector-to-\(\sigma_z\) translation written in this report is inferred from code and indirect tests rather than asserted directly by its own unit test.
\end{itemize}

\section{Potential Ambiguities}
\textbf{Critical:} runtime and synthesis drift layers do not use the same \(\sigma_z\) sign or \(\chi\) scaling.\\
\textbf{Moderate:} higher-order pull/Kerr parameters are falling-factorial coefficients, not ordinary polynomial coefficients.\\
\textbf{Moderate:} the runtime uses \(n_q\) rather than \(\sigma_z/2\), so ground-state energy offsets look different from some textbook forms.\\
\textbf{Low:} the complex-envelope sign \(e^{+i(\omega t+\phi)}\) is repo-specific and must be preserved when adding new pulse families.

\section{Summary}
\texttt{cqed\_sim} is now a coherent reusable library for the qubit-cavity Hamiltonian, frames, pulses, compilation, simulation, standard gates, observables, calibration, tomography, and unitary synthesis. The authoritative physics conventions are: qubit-first tensor ordering; \(|g\rangle=|0\rangle\), \(|e\rangle=|1\rangle\); QuTiP-standard Pauli signs; rad/s and seconds internally; complex envelopes with \(e^{+i(\omega t+\phi)}\); runtime dispersive physics written with \(n_q=|e\rangle\langle e|\); and Lindblad collapse operators as implemented in \texttt{sim/noise.py}. The main caution that remains is the runtime-to-synthesis translation of \(\chi\), which this report now derives explicitly but the code does not yet enforce directly.

\end{document}
